\documentclass[sigconf]{acmart}

% ---- Packages ----
\usepackage{amsmath,amssymb,amsthm}
\usepackage{booktabs}
\usepackage{graphicx}
\usepackage{xcolor}
\usepackage{microtype}
\usepackage{hyperref}
\usepackage{subcaption}

% ---- Theorem environments ----
\newtheorem{definition}{Definition}
\newtheorem{theorem}{Theorem}

% ---- Macros ----
\newcommand{\tva}{\textsf{TVA}}
\newcommand{\tvv}{\textsf{TVV}}
\newcommand{\anchor}{q}
\newcommand{\valq}{\mathrm{Val}_{\anchor}}
\newcommand{\midpoint}{m_V}
\newcommand{\taufill}{\tau_{\mathrm{fill}}}
\newcommand{\tauanchor}{\tau_{\anchor}}

% ================================================================
\begin{document}

\title{Topological Void Validation:\\
Density-Controlled and Role-Aware Evaluation of Predicted Research Voids}

\author{Kris Pan}
\affiliation{\institution{Independent}}
\email{TODO}

% ================================================================
\begin{abstract}
Topological Void Analysis (TVA) identifies candidate research voids in scientific
knowledge spaces as geometric gaps between anchor-conditioned paper clusters. A
natural validation question is whether future papers fill these predicted voids. We
show that raw fill rate---whether any future paper appears near a void's
midpoint---fails as a validation target for three reasons: it is confounded by local
corpus density, limited by anchor observability, and weakly correlated with genuine
epistemic resolution. A hot-zone baseline outperforms TVA on raw fill rate (36\%
vs.\ 28\%), but this apparent advantage disappears under a density-matched baseline
(25\%, TVA/B2 lift $= 1.12\times$, $\mathrm{CI}_{95\%} = [0.85, 1.48]$, not
significant). We further show that anchor eligibility gating exposes a corpus
observability limit: only 5--7\% of future papers are anchor-eligible, so many
unfilled voids should be treated as right-censored. We also demonstrate that a fixed
similarity threshold of 0.82 corresponds to a 35.7\% average null occupancy rate,
ranging from 1\% to 99\% across anchor-density cells; the correct 5\%-FPR threshold
averages 0.841 and is stable across all temporal splits. Under this calibrated
threshold, fill rate drops from $\sim$27\% to $\sim$0.7\%, confirming that most raw
fills are null-zone noise. Finally, role-aware epistemic classification shows that
even anchor-eligible geometric fills rarely correspond to genuine void resolution:
59--76\% are false positives across all groups and splits, with most non-false cases
representing boundary expansion rather than void closure. These findings motivate a
three-layer validation framework---geometric closure, anchor eligibility, and epistemic
role---and show that raw fill rate measures future local activity near predicted voids
but does not by itself validate epistemic closure.
\end{abstract}

\maketitle

% ================================================================
\section{Introduction}
\label{sec:intro}

% TODO: write full intro
% Key sentences:
% - Classical IR: given query, find relevant docs.
%   TVA: given anchor, find missing knowledge.
% - Contributions 1–5

% ================================================================
\section{Background}
\label{sec:background}

\subsection{Topological Void Analysis}
% C1–C4 brief recap, anchor/void/midpoint

\subsection{Temporal Validation in Literature-Based Discovery}
% LBD time-slicing, Recall@K, AP

\subsection{Research Momentum and Dense-Region Confounding}
% Structural holes, hot zones, density bias

% ================================================================
\section{Experimental Setup}
\label{sec:setup}

\subsection{Corpus and Temporal Splits}
% t3/t4/t5 table

\subsection{TVA Void Generation}
% BGE-M3, 10 anchors, top-30 MMR voids

\subsection{Baselines}

\paragraph{B1 (hot-zone baseline).}
% random pairs from top-300 anchor C1 pool

\paragraph{B2 (density-matched baseline).}
% greedy nearest-density match from 300 candidates

\subsection{Three-Layer Fill Predicate}

Let $\valq = \{P \in \mathrm{Val}_t : \mathrm{sim}(P, \anchor) \ge \tauanchor\}$
and
$P_V^* = \arg\max_{P \in \valq} \mathrm{sim}(P, \midpoint)$.

Geometric closure:
\begin{equation}
G(V) = \mathbf{1}\!\left[\mathrm{sim}(P_V^*, \midpoint) \ge \taufill(\anchor, \rho(V), t)\right]
\end{equation}

Validated epistemic fill:
\begin{equation}
\mathrm{Fill}(V) = G(V) \wedge R(P_V^*, A, B, \anchor)
\end{equation}

where $R = 1$ iff $\mathrm{role}(P_V^*, A, B, \anchor) \in
\{\mathrm{TRUE\text{-}FILL},\ \mathrm{PARTIAL\text{-}FILL}\}$.

\paragraph{Calibrated threshold.}
\begin{equation}
\taufill(\anchor, \rho, t) =
Q_{1-\alpha}\!\bigl(
  \{\max_{P \in \valq} \mathrm{sim}(P, m_0) : m_0 \in \mathrm{Null}(\anchor, \rho, t)\}
\bigr)
\end{equation}
where $\mathrm{Null}(\anchor, \rho, t)$ are density-matched null midpoints ($\alpha = 0.05$).

% ================================================================
\section{Results}
\label{sec:results}

\subsection{Finding 1: Raw Fill Rate Is Density-Confounded}
% Table 1 (legacy) + Table 1b (calibrated)

\subsubsection{Fixed Thresholds Are Not Portable}
% Table: FPR@0.82 mean=35.7%, range 1–99%; tau_calib=0.841 stable

\subsection{Finding 2: Anchor Eligibility Reveals Observability Limits}
% Table 2 anchor exposure

\subsection{Finding 3: Geometric Fill Rarely Equals Epistemic Fill}
% Table 3a role decomposition, Table 3b topological event type

\subsection{Case Studies}
% Case 1 PARTIAL_FILL, Case 2 BOUNDARY_EXPANSION, Case 3 FALSE_POSITIVE

% ================================================================
\section{Discussion}
\label{sec:discussion}

\subsection{Raw Fill Fails for Three Reasons}

\subsection{From Fill Detection to Role Decomposition}

\subsection{Descriptive, Not Evaluative}

\subsection{Dynamic Extension (Future Work)}
% One paragraph: papers as events, D-TVA, outside scope

\subsection{Threats to Validity}

% ================================================================
\section{Related Work}
\label{sec:related}

% 1. LBD retrospective validation
% 2. Novelty / first-story detection
% 3. Structural holes and scientometrics
% 4. Dense-region bias in embeddings
% 5. Claim verification / overclaiming

% ================================================================
\section{Conclusion}
\label{sec:conclusion}

Temporal validation of predicted research voids requires separating geometric
occupancy, density bias, anchor observability, and epistemic role. Raw fill rate
alone is insufficient: it conflates void quality with local corpus density and
research momentum. After density matching, the apparent TVA disadvantage
reverses but is not statistically significant. Under calibrated thresholds, fill
rate drops to $\sim$0.7\%, revealing that most raw fills were null-zone noise.
Role-aware classification confirms that geometric fill substantially overestimates
epistemic resolution. We propose a three-layer validation framework and leave
event-driven dynamic extensions to future work.

% ================================================================
\bibliographystyle{ACM-Reference-Format}
\bibliography{paper2}

\end{document}
